
\section{The Palais--Smale condition in an abstract setting}\label{sec:PS-abstract}
\subsection{Energies on submanifolds of Hilbert spaces}
\label{sec:energies-hilbert}
In this section we provide sufficient criteria on the energy and on the underlying
spaces guaranteeing the validity of the Palais--Smale condition (PS). 
Throughout this section we use the following notation:
We let $\HL$ be a Hilbert space and $\BL$ be a Banach space such that there is a compact embedding $\iota_{\BL} \colon \HL \embeds \BL$.
We suppose that $\EL \colon \OL \to \R$ is an energy functional of class $C^1$
on an open subset $\OL \subset \HL$.
Moreover, we suppose that $\ML \subset \OL$ is a $C^1$-submanifold without boundary.
Summarized, we have the following setup:
\begin{enumerate}
	\Item[(E)]\label{condE}
	{\sc (Embedding)}
	$\ML\subset \OL\subset \HL \xhookrightarrow{\iota_\BL} \BL$ and $\iota_\BL$ is compact.
\end{enumerate}
Since $\ML \subset \HL$ is a $C^1$-submanifold, the tangent bundle $T \ML \subset \ML \times \HL$ is splitting, see \cite[Chapter III]{lang_2002}. More precisely, there is a continuous map $\pr^{\ML} \colon \ML \to \BoundedLinOps(\HL)$ such that for each $x \in \ML$ the operator $\pr^{\ML}_x \ceq \pr^{\ML}(x)$ is a continuous linear projector onto $T_x \ML$. Note that we do \emph{not} mean $\pr^{\ML}_x$ to be the orthonormal projector with respect to the inner product on $\HL$.

Our aim is to find sufficient criteria for the \PaSma-condition of $\widetilde{\EL} \ceq \EL|_{\ML}$.
To this end, we make further assumptions.
This will be eased by introducing the following notion:
We say that a subset $\ControlledSet \subset \ML$ is \emph{\Econtrolled{\EL}} if and only if $\ControlledSet \subset \HL$ is bounded and if $\EL$ is bounded on $\ControlledSet$.
Here are our assumptions:
\begin{enumerate}
	\Item[\rm (P)]\label{condP}
	{\sc (Projector)}
	The map $\pr^{\ML}$ is bounded on \Econtrolled{\EL} subsets of $\ML$.
	\Item[\rm (B)]\label{condB}
	{\sc (Boundedness)}
	Every  \PaSma-sequence for $\widetilde{\EL} = \EL|_{\ML}$ is bounded.
	\Item[\rm (L)]\label{condL}
	{\sc (Limit in $\ML$)}
	Every convergent \PaSma-sequence of $\widetilde{\EL}$ has its limit point in $\ML$.
	\Item[\rm (S)]\label{condS}
	{\sc (Splitting)}
	There are 
	a scalar function $\mu \colon \ML \to \intervaloo{0,\infty}$,
	a continuous bilinear form $\QL \colon \HL \times \HL \to \R$,
	and 
	a map $\RL \colon \ML \to \BL^*$
	such that the differential $\D\EL$ splits as follows:
	\[
		\D\EL_x h = \mu(x) \, \QL(x,h) + \RL_x \, \iota_\BL \, h 
		\quad 
		\text{for all $x \in \ML$ and all $h \in\HL$.}
	\]
	Moreover, we assume the following:
	\begin{itemize}
	\Item[\rm (S1)]\label{condS1}
	For every \Econtrolled{\EL} subset $\ControlledSet \subset \ML$ there are $0 < m \leq M < \infty$ such that 
	\[
		m \leq \mu(x) \leq M \quad \text{for all $x\in\ControlledSet$.}
	\]
	
	\Item[\rm (S2)]\label{condS2}
		$\QL$ is coercive, i.e., there is $c > 0$ such that 
	\[
		\QL(h,h) \ge c \norm{h}_\HL^2
		\quad 
		\text{for all $h \in \HL$.}
	\]
	
	\Item[\rm (S3)]\label{condS3}
	$\RL$ is bounded on \Econtrolled{\EL} subsets of $\ML$.
	\end{itemize}

	\Item[\rm (D)]\label{condD}
	{\sc (Differential in non-tangential directions)}
	There is a map $\SL \colon \ML \to \BL^*$ such that 
	\[
		\D{\EL}_x \pars{\Id_\HL-\pr_x^\ML} \, h = \SL_x \, \iota_\BL \, h
		\quad 
		\text{for all $x \in \ML$ and all $h \in \HL$}, 
	\]
	and $\SL$ is bounded on \Econtrolled{\EL} subsets of $ \ML $.
\end{enumerate}

Condition \ref{condS3} is an abstract way of saying $\RL_x$ is of ``lower order'' and thus of less importance compared to $\QL(x,\cdot)$.
Likewise, \ref{condD} states that $\D \EL_x h$ is of ``lower order'' if $h$ is transversal to $T_x \ML$ in the sense that $h$ lies in the range of the $\pars{\Id_\HL-\pr_x^\ML}$.

\begin{theorem}[\PaSma-condition for $\widetilde{\EL}=\EL|_\ML$]
\label{thm:PS-abstract}
Under the conditions \ref{condE}, \ref{condP}, \ref{condB}, \ref{condL}, \ref{condS}, and \ref{condD} above, 
the restricted
energy $\widetilde{\EL}=\EL|_\ML$ satisfies the \PaSma-condition.
\end{theorem}

\begin{proof}
	Let $(x_k)_k\subset\ML$ be a \PaSma-sequence for $\widetilde{\EL}$.
	By condition \ref{condB}, $\ControlledSet \ceq \braces{x_k}_{k \in \N}$ is a bounded subset of $ \HL $. So,
	passing over to a subsequence, we may therefore assume that $(x_{k})_{k \in \N}$ is a weakly convergent \PaSma-sequence. 
	We are going to prove that $(x_{k})_{k \in \N}$ must also be strongly convergent in $\HL$ to a point $x_\infty$;
	because condition \ref{condL} then implies that $x_\infty$ lies in $\ML$, which proves the theorem.

	To show that $(x_{k})_{k \in \N}$ is strongly convergent, we use an argument inspired by the proof of \cite[Lemma 2.1]{cerami-etal_1984}.
	Note that $\ControlledSet \ceq \braces{x_k : k \in \N}$ is \Econtrolled{\EL} by definition of a \PaSma-sequence. 

	Let $0 < m \leq M < 0$ be the constants from \ref{condS1}, and let $c > 0$ be the coercivity constant from \ref{condS2}.
	For $h_k \ceq x_k - x_\infty$ we get
	\begin{equation}
		c \, m \norm{h_k}_\HL^2
		\leq 
		\mu(x_k) \, Q(h_k,h_k)
		=
		\mu(x_k) \, Q(x_k,h_k) - \mu(x_k) \, Q(x_\infty,h_k)
		.
		\label{eq:ExploitingCoercivity}
	\end{equation}
	From \ref{condS} and \ref{condD}, we obtain
	\[
		\mu(x_k) \, Q(x_k,h_k)
		=
		\D \EL_{x_k} \pr^\ML_{x_k} h_k
		+
		\SL_{x_k} \, \iota_\BL \, h_k
		-
		\RL_{x_k} \, \iota_\BL \, h_k
		.
	\]
	The triangle inequality and condition \ref{condS1} lead to
	\begin{align*}
		c \, m \norm{h_{k}}_\HL^2
		&\leq
		\abs{\D \EL_{x_k} \pr^\ML_{x_k} h_k}
		+
		\pars[\big]{
			\norm{\SL_{x_k}}_{\BL^*}
			+
			\norm{\RL_{x_k}}_{\BL^*} 
		} \norm{\iota_{\BL}h_k}_{\BL}
		+
		M \abs{Q(x_\infty,h_k)}
		.
	\end{align*}
	Because $\VL$ is \Econtrolled{\EL},
	conditions \ref{condD}, and \ref{condS3} imply that 
	\[
		C_\SL \ceq \sup_{k \in \N} \norm{\SL_{x_k}}_{\BL^*} < \infty 
		,
		\qand 
		C_\RL \ceq \sup_{k \in \N} \norm{\RL_{x_k}}_{\BL^*} < \infty 
		.
	\]
	Thus, we obtain
	\begin{equation}
		\label{eq:StruweTrick1}
		c \, m \norm{h_{k}}_\HL^2
		\leq
		\abs{\D \EL_{x_k} \pr^\ML_{x_k} h_k}
		+
		\pars[\big]{ C_\SL + C_\RL }
		\norm{\iota_\BL \, h_k}_{\BL}
		+
		M \abs{Q(x_\infty,h_k)}
		.
	\end{equation}
	The bilinear form $Q$ is continuous, and we have weak convergence $h_k \rightharpoonup 0$ in~$\HL$.
	Hence, the third summand $M\abs{Q(x_\infty,h_k)}$ in \cref{eq:StruweTrick1} converges to $0$ as $k \to \infty$.
	The operator $\iota_\BL$ is compact and $h_k$ converges weakly to $0$ in~$\HL$. 
	Hence, $\iota_\BL \, h_k$ converges strongly to~$0$ in~$\BL$.
	Thus, also the second summand in \cref{eq:StruweTrick1} converges to~$0$.
	Finally, we have to discuss the first summand in \cref{eq:StruweTrick1}.
	By definition of the dual norm and the operator norm $\norm{\pr^\ML_{x_k}}_{\LL(\HL)}$, we have
	\[
		\abs{\D \EL_{x_k} \pr^\ML_{x_k} h_k}
		=
		\abs{\D \widetilde{\EL}_{x_k} \pr^\ML_{x_k} h_k}
		\leq
		\norm{\D \widetilde{\EL}_{x_k}}_{(T_{x_k}\ML)^*}
		\norm{\pr^\ML_{x_k}}_{\LL(\HL)}
		\norm{h_k}_{\HL}
		.
	\]	
	Because $(x_k)_{k \in \N}$ is a \PaSma-sequence for $\widetilde{\EL}$, we have $\norm{\D \widetilde{\EL}_{x_k}}_{(T_{x_k} \ML)^*} \to 0$ as $k \to \infty$. 
	Since $\VL$ is \Econtrolled{\EL} and because of \ref{condP},
	we have $\sup_{k \in \N} \norm{\pr^\ML_{x_k}}_{\LL(\HL)} < \infty$. In addition, the weakly convergent sequence $ (h_k)_{k} $ is bounded.
	This shows that also the first summand in \cref{eq:StruweTrick1} converges to $0$. Since $c \, m > 0$, we infer that $h_k \to 0$ strongly in $\HL$, thus $x_k \to x_\infty$ as $k \to \infty$.
\end{proof}

If $\ML= \OL \subset \HL$,
then assumptions \ref{condP}, and \ref{condD} are automatically satisfied, and we obtain a result similar in flavor to \cite[Chapter~II, Proposition 2.2]{struwe_2008a}. 
(Note that the Lax--Milgram theorem implies that coercivity of $\QL$ means that the operator $\HL \to \HL^*$, $x \mapsto \QL(x,\cdot)$ is continuously invertible.)

\begin{corollary}[\PaSma-condition for $\EL$ on $\OL$]
\label{cor:PS-open-set}
Let $\ML=\OL \subset \HL$ be an open set and let $\EL\in C^1(\OL)$. 
Suppose that \ref{condE}, \ref{condB}, \ref{condL},  and  \ref{condS} are satisfied.
Then $\EL$ satisfies the \PaSma-condition.
\end{corollary}

In view of applications in geometric knot theory we consider linear 
combinations of a non-negative energy functional with lower order
terms. 
For this, we extend the abstract setting \ref{condE} by another Banach space $\CL$, into which the Hilbert space $\HL$ embeds compactly, i.e., there is a compact linear embedding $\iota_\CL \colon \HL \to \CL$.
Moreover, we complement condition \ref{condP} on the projection $\pr_x^\ML \colon \HL\to T_x\ML$
by the following condition:

\begin{enumerate}
	\Item[\rm (P*)]\label{condP*}
	There is an operator-valued map $\Pr \colon \ML \to \BoundedLinOps(\CL)$
	such that 
	\[
		\iota_\CL \, \pr^{\ML}_x = \PR_x \, \iota_\CL
		\quad 
		\text{for all $x \in \ML$,}
	\]
	and the map $\Pr$ is bounded on \Econtrolled{\EL} subsets of $\ML$.
\end{enumerate}

In addition to the original energy $\EL\in C^1(\OL)$ with its restriction $\widetilde{\EL}=\EL|_\ML$ to the submanifold $\ML\subset\OL\subset\HL \embeds \BL$, 
we
consider another functional $\FL \in C^1(\UL)$ on an open subset $\UL\subset\CL$ with $\iota_\CL(\OL)\subset\UL$.
For some fixed $\theta>0$ we investigate the  \emph{total energy} $\TL \ceq \EL + \theta \pars{\FL \circ \iota_\CL} \in C^1(\OL)$. 
Note that, in general, a \PaSma-sequence  for the restriction $\widetilde{\TL} \ceq \TL|_\ML$, need not be a \PaSma-sequence for the individual
functionals $\widetilde{\EL}$ or $\widetilde{\FL} \ceq (\FL \circ \iota_\CL) |_\ML$, even if
both functionals are non-negative, since their differentials
may not tend to zero. To formulate a clear-cut result,
we replace assumptions \ref{condL} and
\ref{condB} by the following separate  
assumptions on the energies: 
\begin{enumerate}
	\Item[\rm (B*)]\label{condB*}
	For every $K \in \R$ the sublevel set $\braces[\big]{ x \in \ML : \EL(x) \leq K}$ of $\widetilde{\EL} = \EL_\ML$ is bounded in $\HL$.
	\Item[\rm (L*)]\label{condL*}
	Whenever a sequence $(x_k)_k\subset\ML$ is weakly convergent in $\HL$ to some point $x_\infty$ and $ \sup_{k\in\N}\FL(x_k)<\infty $,
	then $x_\infty \in \ML$.	
\end{enumerate}

\begin{corollary}[(PS)-condition for the total energy $\widetilde{\TL}=\TL|_\ML$]
\label{cor:PS-total-energy}
	Let $\iota_\CL \colon \HL \embeds \CL$ be a compact embedding.
	Let $\OL \subset \HL$ and $\UL \subset \CL$ be open subsets such that $\iota_\CL(\OL) \subset \UL$, and let $\EL \in C^1(\OL)$ and $\FL \in C^1(\UL)$ be bounded from below on $\ML$.
	Suppose that conditions \ref{condE},
	\ref{condP}, \ref{condP*},
	\ref{condB*}, \ref{condS}, \ref{condD},
	and \ref{condL*} are satisfied.
	Then for every $\theta>0$ the restriction  $\widetilde{\TL}=\TL|_\ML$ of the total energy $\TL = \EL + \theta \pars{\FL\circ\iota_\CL}$
	satisfies the \PaSma-condition.
\end{corollary}

\begin{proof}
	By adding suitable constants, we may assume that $\widetilde{\EL} \geq 0$ and $\widetilde{\FL} \geq 0$.
	Let $(x_k)_k\subset\ML$ be a \PaSma-sequence for $\widetilde{\TL}$.
	By definition, there is a $K \in \R$ such that $\widetilde{\TL}(x_k) \leq K$ for all $k \in \N$.
	Now $\widetilde{\FL} \geq 0$ implies
	\begin{equation}
		\label{eq:energy-bound}
		\EL(x_k)
		=
		\widetilde{\EL}(x_k)
		\le
		\widetilde{\EL}(x_k) + \theta \, \widetilde{\FL}(x_k)=\widetilde{\TL}(x_k)\le K\quad\Foa k\in\N
		.
	\end{equation}
	By \ref{condB*}, $(x_k)_{k \in \N}$ is bounded in $\HL$. Hence, it contains a weakly convergent subsequence with limit point $x_\infty \in \HL$. 
	As in the proof of \cref{thm:PS-abstract}, we pass over to this subsequence and denote it by $(x_k)_{k \in \N}$. 
	And as in that proof, we aim to show that the $x_k$ converge strongly to $x_\infty$ in $\HL$.
	Once this is established, condition \ref{condL*} implies $x_\infty \in \ML$
	since $(x_{k})_{k \in \N}$ is bounded under~$\FL$;
	this follows from $\widetilde{\EL} \geq 0$:
	\begin{equation*}
		\FL(x_k)
		=
		\widetilde{\FL}(x_k)
		\le
		\tfrac{1}{\theta}
		\pars[\big]{ \widetilde{\EL}(x_k) + \theta \, \widetilde{\FL}(x_k)}
		=
		\tfrac{1}{\theta} \, \widetilde{\TL}(x_k)
		\le 
		\tfrac{K}{\theta}
		\quad\Foa k\in\N.
	\end{equation*}
	To show strong convergence of $(x_k)_{k \in \N}$, we first recall from \cref{eq:energy-bound} that $\ControlledSet \ceq \braces{x_k}_{k \in \N}$  is \Econtrolled{\EL}.
	Precisely as in the proof of \cref{thm:PS-abstract}, we derive inequality \cref{eq:StruweTrick1} for $h_k \ceq x_k - x_\infty$ and show that the second and third summand in \cref{eq:StruweTrick1} converge to~$0$. Alas, we need a new argument for showing that the first summand $\abs{\D \EL_{x_k} \pr^\ML_{x_k} h_k}$ converges to~$0$ as $ k\to \infty $.
	This time, we decompose it and use condition \ref{condP*} as follows:
	\[
		\D \EL_{x_k} \pr^\ML_{x_k} h_k
		=
		\D \TL_{x_k} \pr^\ML_{x_k} h_k
		-
		\theta \, \D \FL_{x_k} \, \iota_\CL \, \pr^\ML_{x_k} h_k
		=
		\D \widetilde{\TL}_{x_k} \pr^\ML_{x_k} h_k
		-
		\theta \, \D \FL_{(\iota_\CL x_k)} \, \PR_{x_k}\, \iota_\CL \,h_k
		.
	\]
	This leads us to
	\[
		\abs{\D \EL_{x_k} \pr^\ML_{x_k} h_k}
		\leq 
		\norm{\D \widetilde{\TL}_{x_k}}_{(T_{x_k}\ML)^*}
		\norm{\pr_{x_k}^\ML}_{\LL(\HL)}
		\norm{h_k}_{\HL}
		+
		\theta \norm{\D \FL_{(\iota_\CL x_k)}}_{\CL^*}
		\norm{\PR_{x_k}}_{\LL(\CL)}
		\norm{\iota_\CL \,h_k}_{\CL}
		.
	\]
	Since $\iota_\CL$ is compact, we have $\iota_\CL \, x_k \to \iota_\CL \, x_\infty$ as $k \to \infty$. 
	By condition~\ref{condL*}, we have $x_\infty \in \ML \subset \OL$
	and $\iota_\CL x_\infty \in \iota_\CL(\OL) \subset \UL$.
	Since $x \mapsto \D \FL_x$ is continuous on $ \UL $, and by condition~\ref{condP}, and because of $\iota_\CL \, x_k \to \iota_\CL \, x_\infty$, we have
	\[
		\smash{\sup_{k \in \N}} \norm{\D \FL_{(\iota_\CL x_k)}} < \infty
		,
		\quad 
		\sup_{k \in \N} \norm{\pr_{x_k}^{\ML}}  < \infty
		,
		\qand 
		\smash{\lim_{k \to \N}} \norm{\iota_\CL \, h_k}_{\CL} = 0
		.
	\]
	The sequence $(x_k)_{k\in\N}$ is a \PaSma-sequence for $\smash{\widetilde{\TL}}$, and because of \ref{condP*}, and $h_k \rightharpoonup 0$, we have 
	\[
		\smash{\lim_{k \to \infty}}
		\norm{\D \widetilde{\TL}_{x_k}}_{(T_{x_k}\ML)^*}
		=
		0
		,
		\quad 
		\sup_{k \in \N} \norm{\PR_{x_k}}_{\LL(\CL)}  < \infty
		,
		\qand
		\smash{\sup_{k \in \N}} \norm{h_k}_{\HL}  < \infty
		.
	\]
	All this together implies that $\smash{\abs{\D \widetilde{\EL}_{x_k} \pr^\ML_{x_k} h_k}} \to 0$ as $k \to \infty$, which completes the proof.
\end{proof}

\subsection{Gradient flows on Riemannian manifolds}
\label{sec:grad_flows}
Let $(\ML,g)$ be a Riemannian manifold without boundary 
of class $C^2$ over a Hilbert 
space $\HL$ with  Riemannian metric $g$ given fibrewise as 
$g_x(\cdot,\cdot)$ 
inducing a norm $\|\cdot\|_{g_x} \ceq \sqrt{g_x(\cdot,\cdot)}$ on 
the tangent space $T_x\ML$. 
For two points $x,y$ 
in the same connected component of $\ML$ we consider
the  set of connecting $C^1$-paths,
\begin{equation}\label{eq:path-space}
	\PS(x,y) 
	\ceq \braces[\big]{c\in C^1([0,1],\ML) : c(0)=x,c(1)=y},
\end{equation}
and define the \emph{path metric} (or \emph{Palais distance})
\begin{equation}\label{eq:path-metric}
\dist_{\ML}(x,y) \ceq \inf_{c\,\in\PS(x,y)}\int_0^1
\|\dot{c}(t)\|_{g_{c(t)}}\d t,
\end{equation}
where $\dot{c}(t) \ceq \D c_t(1)\in T_{c(t)}\ML$
for the differential $\D c_t:T_t[0,1]\simeq\R\to
T_{c(t)}\ML$ of the $C^1$-map $c:[0,1]\to\ML$ (cf.{} \cite[Sections 6 \& 9]{palais_1963}).
Equipped with the path metric $\dist_{\ML}$, every connected
component
of 
$\ML$ is a metric space. This metric is consistent with the
topology of $\ML$ as shown in \cite[Section 9]{palais_1963}
and in the more general setting of Finsler manifolds in
\cite[Theorem 3.3]{palais_1966}. Therefore, every connected
component $\KL$ of $\ML$ is closed with respect to $\dist_{\ML}$;
see, e.g., \cite[Remark after Theorem 25.2]{munkres_2000}.

On this manifold we consider energy functionals $\EL:\ML\to\R$ 
of class $C^{1,1}_\loc$ whose gradient $\nabla \EL$ is
given by $\D\EL_x(h)=g_x(\nabla\EL(x),h)$ for $x\in\ML$ and $h\in
T_x\ML$. We assume that
\begin{enumerate}
	\Item[\rm (e)]\label{lower-bound} 
	({\sc Energy bounded below})
	There is a number $e\in\R$ such that
	$\EL(x)\ge e$ for all $x\in\ML$.
	\Item[\rm (L**)]\label{condL**} 
	({\sc limit in $\ML$}) 
	Every Cauchy sequence
	$(x_k)_{k \in \N}$ in $\ML$ that is bounded under $\EL$
	converges to some 
	$x_\infty\in\ML$ as $k\to\infty$.
\end{enumerate}
Long-time existence for the gradient flow
and subconvergence to critical points can
now be proved by slightly generalizing a classic argument 
that can be found, e.g., in \cite[Theorem 9.1.6]{palais-terng_1988} 
for Riemannian manifolds, in \cite[Theorem 1.1]{browder_1965} 
for Banach manifolds modelled over a uniformly convex Banach space, or 
in \cite[Theorem 5.4]{palais_1966} for general Finsler manifolds.

\begin{theorem}[Gradient flow]
	\label[theorem]{thm:gradflow_subconvergence}
	Let $\EL\in C^{1,1}_\loc(\ML)$ satisfy the Palais--Smale condition
	{\rm (PS)} and conditions \ref{lower-bound} and \ref{condL**}.
	There exists a function $\phi:[0,\infty)\times\ML\to\ML$, such
	that $\phi(\cdot,x)\in C^1([0,\infty),\ML)$ for every $x\in\ML$,
	and
	\begin{equation}\label{eq:gradient-flow}
	\begin{cases}
	\frac{\d}{\d t}\phi(t,x)& =-\nabla\EL(\phi(t,x)) \Foa t\in [0,\infty),\\
		\phi(0,x)&=x.  
	\end{cases}
	\end{equation}
The function $t\mapsto\EL(\phi(t,x))$ is strictly decreasing 
on $[0,\infty)$ unless
$x$ is
	$\EL$-critical.
For $x\in\ML$ and $t\to\infty$ the solution $\phi(t,x)$ subconverges  
to a critical point of $\EL$, which lies in the same
connected component of $\ML$ as the initial point $x$.
In particular, in each connected component
of $\ML$ the energy $\EL$ possesses at least one
critical point. 
\end{theorem}

\begin{proof}
	For $x\in\ML$ we obtain
	the function $\phi$ as the solution to the 
	initial value problem
	\begin{equation} 
		\begin{aligned}
			\label{eq:cauchy_problem}
		\textstyle
		\frac{\d}{\d t}\xi(t) &=-\nabla \EL(\xi(t)), &\hspace{10pt}\xi(0)&=x.
		\end{aligned}
	\end{equation}
	Such a solution exists and is unique on an 
	interval $(t^-,t^{+})$ containing $t=0$
	by the Picard--Lindelöf Theorem 
	(see, e.g., \cite[Chapter IV.2]{lang_2002}). 
	Here $t^{+}\in (0,\infty]$ 
	denotes the maximal positive existence time, 
	and for $t\in (t^-,t^+)$ we compute by means of 
	\cref{eq:cauchy_problem}
	\begin{align}
	\textstyle
	\frac{\d}{\d t}\EL(\xi(t)) &
=\D\EL_{\xi(t)}(\dot{\xi}(t))= g_{\xi(t)}\big(\nabla \EL(\xi(t)),
	\dot{\xi}(t)\big)\label{eq:decreasing-energy}\\
		&\overset{\cref{eq:cauchy_problem}}{=}  
- g_{\xi(t)}\big(\nabla \EL(\xi(t)),\nabla \EL(\xi(t))\big)
=-\|\nabla\EL(\xi(t))\|^2_{g_{\xi(t)}}\le 0.\notag
	\end{align}
Hence, $\phi(\cdot,x) \ceq \xi(\cdot)$ 
solves 
\cref{eq:gradient-flow} on the positive maximal time interval $[0,t^+)$, 
and the energy is strictly decreasing along the flow unless $\D\EL_x=0$,
or equivalently, $\nabla\EL(x)=0$, in which case $\xi(t)=x$ and therefore
$\EL(\xi(t))=\EL(x)$ for all $t\in
[0,t^+)$ by uniqeness of the solution to \cref{eq:cauchy_problem}.

To show that $t^+=\infty$ we assume to the contrary that $t^+<\infty$,
and show below
that in that case the one-sided limit
$\xi(t^+) \ceq \lim_{t\nearrow t^+}\xi(t)$ exists, so that 
we could find a solution of \cref{eq:cauchy_problem} with $x$ replaced
by $\xi(t^+)$ on an
open
 time interval containing $t^+$, which contradicts the 
maximality of $t^+$ by uniqueness of the solution.

Towards the existence of the limit point
$\xi(t^+)\in\ML$ for some finite maximal time $t^+\in
	(0,\infty)$
consider 
a strictly increasing
sequence $(t_k)_k\subset[0,t^+)$ with $t_k\to t^+$ as $k\to\infty$,
and define 
the sequence of points
$x_k=\xi(t_k)\in\ML$. We first show that $x_k$ is a Cauchy 
sequence in $\ML$, so that by condition
\ref{condL**} we find $x_\infty\in\ML$ such that $x_k\to x_\infty$ as
$k\to\infty$, since we know  that
$\EL(x_k)\le\EL(x)$ for all $k\in\N$ by \cref{eq:decreasing-energy}. Going back
to an arbitrary sequence $s_m\nearrow t^+$ we will then show that also 
$y_m \ceq \xi(s_m)\to x_\infty$ as $m\to\infty$, so that $\xi(t^+)$ indeed exists.

We have by \cref{eq:decreasing-energy} 
\begin{align}
	\EL(\xi(t))-\EL(\xi(s))\overset{\cref{eq:decreasing-energy}}{=}
	\textstyle
	-\int_s^t \norm{ \nabla \EL(\xi(\tau))}^2_{g_{\xi(\tau)}}\d \tau
	\Fo 0\le s< t<t^+,\label{eq:energy-identity}
\end{align}
which yields the following by means of \ref{lower-bound}:
\begin{align*}
\textstyle
\int_{s}^{t} \norm{\nabla \EL(\xi(\tau))}_{g_{\xi(\tau)}}
\d \tau 
&\textstyle
	\le 
	\pars{t-s}^{1/2} 
	\pars[\big]{ \int_s^t
		\norm{\nabla \EL(\xi(\tau))}^2_{g_{\xi(\tau)}}\d \tau 
	}^{1/2}\\
&\overset{\ref{lower-bound}}{\le}(t-s)^{\frac{1}{2}}(\EL(x)-e)^{\frac{1}{2}}
\quad\Foa 0\le s<t<t^+.
	\end{align*}
	Hence, by definition \cref{eq:path-metric} of the path metric,
	and by the differential equation \cref{eq:cauchy_problem},
	\begin{align}
		\dist_\ML(x_j,x_k)&\overset{\cref{eq:path-metric}}{\le}
		\textstyle
		\int_{t_j}^{t_k}\|\dot{\xi}(\tau)\|_{g_{\xi(\tau)}}\d
		\tau \overset{\cref{eq:cauchy_problem}}{=}
		\int_{t_j}^{t_k}\|\nabla \EL(\xi(\tau)) \|_{g_{\xi(\tau)}}
		\d 
		\tau\notag\\
		&\le (t_k-t_j)^{\frac{1}{2}}(\EL(x)-e)
		\longrightarrow 0 \As j,k\to\infty;\label{eq:cauchy-estimate}
	\end{align}
hence $(x_k)_k$ is indeed a Cauchy sequence so that $x_k\to x_\infty$
for some $x_\infty\in\ML$. For any other sequence $s_m\nearrow t^+$ repeat
the argument in \cref{eq:cauchy-estimate} with $y_m \ceq \xi(s_m)$ instead of
$x_j$ to find
\begin{align*}
	\dist_\ML(y_m,x_\infty)
	=
	\lim_{k\to\infty} \dist_\ML(y_m,x_k)
	&\le
	\lim_{k\to\infty} \pars{t_k-s_m}^{1/2} \pars{\EL(x)-e}^{1/2}
	\\
	& = 
	\pars{t^+-s_m}^{1/2} \pars{\EL(x)-e}^{1/2}
	\to 
	0
	\As m\to\infty,
\end{align*}
so that the one-sided limit $\xi(t^+)$ indeed exists.
Therefore, 
	the solution $\phi(\cdot,x)$ exists on $[0,\infty)$.

Combining the energy identity \cref{eq:energy-identity} for 
$t \ceq t_k\nearrow\infty$ and 
$s \ceq 0$ with condition \ref{lower-bound}
we obtain
$$
\textstyle
\int_0^\infty\|\nabla\EL(\xi(\tau))\|^2_{g_{\xi(\tau)}}\d \tau=
\lim_{k\to\infty}\int_0^{t_k}\|\nabla\EL(\xi(\tau))\|^2_{g_{\xi(\tau)}}\d \tau
\le \EL(x)-e.
$$
Therefore, there exists a sequence $s_l\nearrow\infty$ such that
for the points $z_l \ceq \xi(s_l)\in\ML$ one has
$\|\nabla\EL(z_l)\|_{g_{z_l}}\to 0$ 
as $l\to\infty$. Together with 
$\EL(z_l)\le\EL(x)$ for all $l\in\N$ this shows that
$(z_l)_l\subset\ML$ is a (PS)-sequence, which by means of the (PS)-condition
implies
that there is a subsequence $(z_{l_i})_i\subset (z_l)_l$ 
converging to some point
$z_\infty\in\ML$ as $i\to\infty$. Since $\EL$ is of class $C^1$ we find
$$
\|\D\EL_{z_\infty}\|_{(T_{z_\infty}\ML)^*}=\lim_{i\to\infty}
\|\D\EL_{z_{l_i}}\|_{(T_{z_{l_i}}\ML)^*}=\lim_{i\to\infty}
\|\nabla\EL(z_{l_i})\|_{g_{z_{l_i}}}=0,
$$
that is, $z_\infty$ is $\EL$-critical. 
All points $z_{l_i}$ are connected to the initial point $x$ by
the continuous path $\xi|_{[0,s_{l_i}]}$, so they lie in the same
connected component $\KL$ of $\ML$ as $x$. Since $\KL$ is closed,
one also has $z_\infty\in\KL$.
\end{proof}
Palais considers in \cite[Section 2]{palais_1979} a group $G$ of
isometries of a Riemannian manifold and a $G$-invariant energy $\EL:\ML\to
\R$ to provide an exemplary simple setting where his
principle of symmetric criticality holds true. It turns out that in this
situation the gradient flow \cref{eq:cauchy_problem} is symmetry
preserving.
\begin{proposition}[Symmetry-preserving gradient flow]
\label{prop:symmetry-preserving-grad-flow}
Let $G$ be a group of $C^1$-isometries of the Riemannian manifold $(\ML,g)$. 
Suppose
that $\EL\in C^{1,1}_\textnormal{loc}(\ML)$ is invariant under $G$, i.e.,
$\EL=\EL\circ\G$ for all $\G\in G$, and satisfies the assumptions of
\Cref{thm:gradflow_subconvergence}. Then the gradient flow 
\cref{eq:cauchy_problem} preserves the $G$-symmetry, that is,
for every initial point $x\in\ML$ that satisfies $\G x=x$ for all $\G\in G$, the solution
$\xi$ of \cref{eq:cauchy_problem}  satisfies $\G \xi(t)=\xi(t)$ for all
$\G \in G$ and $t\in [0,\infty)$. In addition, $\xi$ subconverges to a 
$G$-symmetric critical point $x_\infty$, i.e., with $\G x_\infty = x_\infty$
for all $\G\in G$.
In particular, in any connected component $\KL$ of $\ML$ that contains
a $G$-symmetric point $x$ there is a $G$-symmetric $\EL$-critical point.
\end{proposition}
\begin{proof}
Palais showed in \cite[Section 2]{palais_1979} that 
$\nabla\EL(\G x)=
\D\G_x(\nabla\EL(x))$ for all $\G\in G$ and $x\in\ML$. 
This implies for any $G$-symmetric
initial point $x\in\ML$ and any fixed group element $\G\in G$ that
the solution $\xi$ of \cref{eq:cauchy_problem}
generates the
function $\varXi(t) \ceq \G\xi(t)$, which satisfies $\varXi(0)=\G\xi(0)=
\G x=x$ and
$$
\dot{\varXi}(t)=\D\G_{\xi(t)}(\dot\xi(t))\overset{\cref{eq:cauchy_problem}}{=}
-\D\G_{\xi(t)}\nabla\EL(\xi(t))=-\nabla\EL\G\xi(t)=-\nabla\EL\varXi(t)
$$
for all $t\in [0,\infty)$. 
In other words, $\varXi$ solves the same initial
value problem as $\xi$ does, so that $\varXi=\xi$ by uniqueness of the solution.
Since $\G\in G$ was chosen arbitrarily, we have shown that \cref{eq:cauchy_problem} is 
symmetry-preserving. If $x_\infty\in\ML$ is an $\EL$-critical limit point of 
$\xi$ then there is a sequence 
$t_i\to\infty$ such that 
$$
\G x_\infty=\G \lim_{i\to\infty}\xi(t_i)=\lim_{i\to\infty}\G\xi(t_i)=
\lim_{i\to\infty}\xi(t_i)=x_\infty,
$$
so $x_\infty $ is $G$-symmetric as well. Finally, to find a $G$-symmetric
$\EL$-critical point in a component $\KL$ of $\ML$ containing a $G$-symmetric
point $x\in\KL$, we simply start the gradient flow \cref{eq:cauchy_problem}
in that point $x$. This flow subconverges to a 
$G$-symmetric $\EL$-critical point
in that same component $\KL$.
\end{proof}
\begin{corollary}[Critical points of lower energy]
{\cite[Corollary 9.1.8]{palais-terng_1988}}
	\label[corollary]{koro:critpoint_lowerenergy}
Let $\EL$ satisfy the assumptions of \Cref{thm:gradflow_subconvergence}.
Then, for any point $x\in\ML$ with $\D\EL_x\not=0$ there
exists a critical point $x_{\text{c}}$ of $\EL$ in the connected component
of $\ML$ that contains $x$, such that $\EL(x_{\text{c}})<\EL(x)$.
If $x$ is $G$-symmetric and $\EL$ is $G$-invariant under a group
$G$ of $C^1$-isometries of $(\ML,g)$ then $x_c$ is $G$-symmetric as well.
\end{corollary}
\begin{proof}
	Choose $x_{\text{c}}\in\ML$ as any limit point of $\phi(t,x)$
	as $t\to\infty$, 
	which is contained in the same connected
	component $\KL\subset\ML$ as the initial point $x$ according
	to \Cref{thm:gradflow_subconvergence}. 
	Since $x$ is not a critical point we have $\EL(x_\text{c})<\EL(x)$
	because of the strict energy decay along the flow. The last
	statement follows from \Cref{prop:symmetry-preserving-grad-flow}.
\end{proof}
\begin{corollary}[Energy minimizers]{\cite[Theorem 9.1.9]{palais-terng_1988}}
	\label[corollary]{koro:minimizer_existence}
	Under the assumptions of \Cref{thm:gradflow_subconvergence}
	 the energy $\EL$ attains its infimum on $\ML$, and additionally on 
 every connected component of $\ML$. Moreover, on $\ML$, and also
 in each of its connected components
 there exists a minimizing {\rm (PS)}-sequence. 
 If a connected component
 $\KL$ of $\ML$ contains a $G$-symmetric point $x$ and $\EL$ is $G$-invariant
 under a group $G$ of $C^1$-isometries 
 of $(\ML,g)$, then there is a $G$-symmetric
 element $x_s\in\KL$ with $\EL(x_s)=\inf\{\EL(y):y\in\KL,\,\textnormal{$y$
 is $G$-symmetric}\}$, and a corresponding {\rm (PS)}-sequence of
 $G$-symmetric elements in $\KL$.
\end{corollary}
\begin{proof}
Let $x_k$ be a minimizing sequence for $\EL$ on $\ML$, i.e.,
$\lim_{k\to\infty}\EL(x_k)=\inf(\EL)$. 
Without loss of generality we can assume, that each $x_k$ is a critical point, since otherwise we just replace $x_k$ 
by a critical point of lower energy by \cref{koro:critpoint_lowerenergy}. Then, since $\EL$ satisfies the (PS)-condition, 
there exists a subsequence $(x_{k_i})_i\subset (x_k)_k$ 
converging to a critical point 
$x_{\text{c}}$ as $i\to\infty$, and by continuity of $\EL$ one has
$
\EL(x_{\text{c}})=\lim_{i\to\infty}\EL(x_{k_i})=\inf(\EL).
$
We can repeat the same argument on each connected component of $\ML$ instead
of $\ML$, since each such component is closed.
Finally, if we consider a minimizing sequence of $G$-symmetric
elements $x_k\in\KL$, then the same reasoning is valid due to the last
statement of \cref{koro:critpoint_lowerenergy} and because of the fact
that $G$-symmetry is a closed condition since each $\G\in G$ is a
continuous mapping from $\ML$ onto $\ML$.
\end{proof}


Even if the energy is only in $C^1$, we can use Ekeland's variational
principle to obtain energy-minimizing
Palais--Smale sequences subconverging to
energy minimizers. 

\begin{corollary}[Minimizers]
\label{cor:critical+minimizers}
Let $\EL\in C^1(\ML)$ satisfy the Palais--Smale condition {\rm (PS)} and conditions \ref{lower-bound} and \ref{condL**}.
Let $\KL$ be a connected component of $\ML$.
Then there are a minimizing {\rm (PS)}-sequence and a minimizer of $\EL |_{\KL}$.

If  $\KL$ contains a $G$-symmetric element and if
 $\EL$ is $G$-invariant, then
there exists a (not necessarily $G$-symmetric) {\rm (PS)}-sequence for $\EL|_\KL$ that converges to
a $G$-symmetric element in $\KL$ that realizes $\inf\{\EL(y):y\in\KL,\,\textnormal{$y$
is $G$-symmetric}\}$.
These statements hold true also for $\KL = \ML$.
\end{corollary}
\begin{proof}
We  choose a minimizing sequence $(x_k)_{k \in \N}$ of $\EL |_{\KL}$ such that $\EL(x_k) \le \inf(\EL|_{\KL}) + 1/k^2$.
The function $\EL$ is continuous and the set $\KL$ is closed.
Hence, condition \ref{condL**} implies that the sublevel set 
$
	\mathcal{U} \ceq \braces[\big]{x \in \KL : \EL(x) \leq \inf (\EL|_\KL) + 2}
$ 
equipped with the restriction of $\dist_\ML$ is a complete metric space.
Now we apply Ekeland's variational principle (see, e.g., \cite[Chapter~1,Theorem~5.1]{struwe_2008a}) to find $y_k \in \mathcal{U}$
such that 
$
	\EL(y_k) \le \EL(x_k)
$,
$
	\dist_\ML(x_k,y_k) \le 1/k
$,
and
\begin{equation}
	\label{eq:ekeland-ineq}
	\EL(y_k) \leq \EL(z) + \tfrac{1}{k} \dist_\ML(y_k,z) \quad \Foa z\in\mathcal{U}
	.
\end{equation}
We have $\EL(y_k) \leq \EL(x_k) \leq \inf(\EL|_{\KL}) + 1 < \inf(\EL|_\KL) + 2$.
Because $\EL$ is continuous and because $\KL$ is open, $ y_k $ is contained in the interior of $ \UL $.
So, for every tangent vector $h\in T_{y_k}\ML$ 
there is a $\sigma > 0$ and a $C^1$-path $c \colon \intervaloo{-\sigma,\sigma} \to \mathcal{U}$ satisfying $c(0)=y_k$ and $\dot c(0) = h$.
By continuity, we can take $\sigma$ so small that
\begin{equation}\label{eq:tangent-small}
	\textstyle
	\norm{\dot c(t)}_{g_{c(t)}}\le 2\norm{h}_{g_{c(0)}} = 2 \norm{h}_{T_{y_k}\ML}
\quad\Foa t\in (-\sigma,\sigma).
\end{equation}
Applying \cref{eq:ekeland-ineq} to the point
$z \ceq c(t)$ for any $t\in (0,\sigma)$
we obtain
\begin{equation}\label{eq:PS-weg}
\textstyle
	\EL(y_k)-\EL(c(t))
	\overset{\cref{eq:ekeland-ineq}}{\le}
	\frac1{k}\dist_\ML(y_k,c(t)) \le \frac1{k}\int_0^t
	\norm{c'(\tau)}_{g_{c(\tau)}}\d \tau\le \frac2{k}\norm{h}_{T_{y_k}\ML}\cdot t
	.
\end{equation}
Dividing by $t>0$ and taking the limit $t\searrow 0$,
we arrive at
$\D \EL_{y_k} h \le \frac{2}{k} \norm{h}_{T_{y_k}\ML}$, hence
\[ 
	\norm{\D\EL_{y_k}}_{(T_{y_k}\ML)^*} \le \tfrac{2}{k}\to 0
	\quad  \text{as} \quad  k \to\infty
	.
\] 
Thus, $(y_k)_k$ is a minimizing (PS)-sequence for $\EL|_{\KL}$.
In particular, it is a (PS)-sequence for $\EL$.
Because $\EL$ satisfies the (PS)-condition, the sequence subconverges to some $y_\infty \in \ML$: $y_{k_i}\to y_\infty$ as $i\to\infty$.
Since $\mathcal{U}$ is closed, we have $y_\infty \in \mathcal{U} \subset \KL$, hence
\[
	\inf(\EL|_{\KL})
	\leq 
	\EL(y_\infty)
	=
	\lim_{i\to\infty} \EL(y_{k_i})
	\le
	\lim_{i\to\infty} \EL(x_{k_i})
	=
	\inf(\EL|_{\KL})
	. 
\]
Thus, $y_\infty$ is the desired minimizer.

If $\EL$ is $G$-invariant,
we can start with a minimizing sequence $(x_k)_k\subset\KL$ of $G$-symmetric
elements, and then follow the same line of arguments to show
that a subsequence $(x_{k_i})_i\subset (x_k)_k$ 
converges to some $y_\infty\in \KL$ as $i\to\infty$. Because $G$-symmetry is a closed condition, we see that the limit point $y_\infty$ is $G$-symmetric, too. Finally, we notice that the above argument goes through, too, if we replace $\KL$ by $\ML$.
\end{proof}